% =============================================================================
% 5 Detaillierte Ausarbeitung des gewählten Konzepts (05_detaillierte_ausarbeitung.tex)
% =============================================================================
\section{Detaillierte Ausarbeitung des gewählten Konzepts}
\label{chap:ausarbeitung}

Basierend auf der Evaluierung in Kapitel 4 wird im Folgenden die hybride ohmsch-induktive Lastsimulation detailliert ausgearbeitet. Die Auslegung erfolgt exemplarisch für das anspruchsvollste Prüfmodul (Typ 3er, 630\,A Dreiphasen-Wechselstrom). Ziel dieses Kapitels ist die physikalisch-mathematische Dimensionierung der Hardwarekomponenten sowie die Überführung des theoretischen Konstrukts in eine reale Automatisierungsarchitektur mittels Siemens ET200 SP.

% -----------------------------------------------------------------------------
% 5.1 Hardware-Konzeption und Dimensionierung
% -----------------------------------------------------------------------------
\subsection{Hardware-Konzeption und Dimensionierung}
\label{sec:hardware_konzeption}

Die hybride R-L-Kombination erfordert eine exakte mechanische und elektrische Abstimmung zwischen dem ungesteuerten Festwiderstand und der regelbaren Tauchkerndrossel. Die Dimensionierung muss gewährleisten, dass das System auch bei einer achtstündigen Prüfdauer unter Volllast thermisch und magnetisch stabil bleibt.
% -----------------------------------------------------------------------------
% 5.1.1 Auslegung des ohmschen Hauptpfades

\subsubsection{Auslegung des ohmschen Hauptpfades und thermodynamische Modellierung}
\label{sec:auslegung_widerstand}

Der Festwiderstand hat die Aufgabe, den dominierenden Wirkstromanteil ($I_R \approx 500\,\text{A}$) zu führen, um den normativ geforderten Leistungsfaktor ($\cos\varphi \approx 1$) für den Bauartnachweis zu sichern. Als Widerstandsmaterial wird eine hochlegierte Chrom-Nickel-Stahllegierung (Werkstoffnummer 1.4310, X10CrNi18-8) spezifiziert. Diese Legierung zeichnet sich durch eine hohe Zunderbeständigkeit bis zu $800\,^\circ\text{C}$ und einen moderaten spezifischen elektrischen Widerstand aus.

Die abzuführende thermische Verlustleistung $P_v$ im Widerstand bei Nennstrom berechnet sich nach dem Joule'schen Gesetz. Bei einem Kaltwiderstand von $R_{20} = 8,1\,\text{m}\Omega$ ergibt sich:
\begin{equation}
    P_v = I_R^2 \cdot R_{20} \approx (500\,\text{A})^2 \cdot 8,1\,\text{m}\Omega = 2.025\,\text{W}
\end{equation}
Da in einem dreiphasigen System über $6\,\text{kW}$ Abwärme pro Modul entstehen, ist das thermische Management von zentraler Bedeutung. Das transiente Erwärmungsverhalten des Festwiderstands wird maßgeblich durch seine Bauform bestimmt. Um die wirksame Kühloberfläche $A_{eff}$ zu maximieren, wird das Widerstandsmaterial nicht als kompakter Block, sondern als gefaltetes Mäander-Gitter ausgeführt. Diese Geometrie begünstigt den konvektiven Wärmeübergang (Kamineffekt), da die erwärmte Luft ungehindert vertikal aufsteigen kann und kühlere Umgebungsluft von unten nachzieht.

Die Erwärmung folgt der homogenen Wärmeleitungsgleichung, welche die zugeführte elektrische Energie der gespeicherten und der abgeführten Wärmeenergie gleichsetzt:
\begin{equation}
    \vartheta(t) = \vartheta_{Umg} + \frac{P_v}{A_{eff} \cdot \alpha_w} \cdot \left( 1 - e^{-\frac{t}{\tau_{th}}} \right)
\end{equation}
Die thermische Zeitkonstante $\tau_{th} = (m \cdot c_p) / (A_{eff} \cdot \alpha_w)$ ist für die Regelungstechnik von enormem Vorteil. Durch die massive Bauweise des Edelstahlgitters (große Masse $m$ und spezifische Wärmekapazität $c_p$) liegt $\tau_{th}$ im Bereich von mehreren Minuten. Dies führt zu einem stark geglätteten, asymptotischen Temperaturanstieg. Kurzzeitige Netzschwankungen am Haupttransformator wirken sich dadurch nicht unmittelbar auf den Widerstandswert aus, was der ET200 SP ausreichend Zeit für ausgleichende Stellbewegungen an der Drossel gewährt.
% -----------------------------------------------------------------------------
% 5.1.2 Magnetkreisberechnung
\subsubsection{Magnetkreisberechnung und Sättigungsprävention der Tauchkerndrossel}
\label{sec:auslegung_drossel}

Die parallel geschaltete Tauchkerndrossel kompensiert den thermischen Widerstandsanstieg durch die Aufnahme eines steuerbaren Blindstroms $I_L$. In der komplexen Wechselstromrechnung wird der Gesamtstrom $I_{ges}$ durch die vektorielle Addition von Wirk- und Blindstrom gebildet. Da der Wirkstrom auf der reellen Achse und der Blindstrom der Induktivität um $-90^\circ$ phasenverschoben auf der imaginären Achse liegt, ergibt sich der Betrag zu:

\begin{equation}
    I_{ges} = \sqrt{I_R^2 + I_L^2}
\end{equation}
Sinkt der Wirkstrom durch die Eigenerwärmung des Stahlgitters auf beispielsweise $440\,\text{A}$ ab, muss die Steuerung die Drossel so verfahren, dass ein Blindstrom von $I_{L,max} = \sqrt{630^2 - 440^2} \approx 451\,\text{A}$ fließt, um den Prüfstrom exakt bei 630\,A zu halten.

Die größte physikalische Herausforderung bei Hochstromdrosseln ist die Vermeidung der magnetischen Sättigung. Die Magnetisierungskennlinie (B-H-Kurve) von kornorientiertem Elektroblech verläuft nur bis zu einer magnetischen Flussdichte von $B_{sat} \approx 1,5\,\text{T}$ linear. Wird dieser Punkt überschritten, bricht die Induktivität zusammen, was zu einer massiven, nicht-sinusförmigen Verzerrungsblindleistung führt. Um dies zu verhindern, wird der Eisenquerschnitt großzügig auf $A_{Fe} = 30\,\text{cm}^2$ ($0,003\,\text{m}^2$) bei $N=8$ Windungen dimensioniert, wodurch die maximale Flussdichte auf unkritische $\approx 0,47\,\text{T}$ limitiert wird.

Nach dem Hopkinson'schen Gesetz (magnetischer Maschensatz) setzt sich der magnetische Gesamtwiderstand $R_{m,ges}$ der Tauchkerndrossel aus dem Eisenweg und dem verstellbaren Luftspalt $\delta(x)$ zusammen:
\begin{equation}
    R_{m,ges} = \frac{l_{Fe}}{\mu_0 \cdot \mu_r \cdot A_{Fe}} + \frac{\delta(x)}{\mu_0 \cdot A_{\delta}}
\end{equation}
Da die relative Permeabilität des Elektroblechs ($\mu_r \gg 1000$) extrem hoch ist und das Eisen weit unterhalb der Sättigung betrieben wird, stellt der Kern fast keinen magnetischen Widerstand dar. Zur Berücksichtigung der physikalischen Streuung (Feldaufweitung am Luftspalt) wird ein um $10\,\%$ vergrößerter wirksamer Spaltquerschnitt $A_\delta = 1,1 \cdot A_{Fe} = 0,0033\,\text{m}^2$ angesetzt. 
Die resultierende stufenlose Induktivität in Abhängigkeit der Kernposition $\delta(x)$ lautet in sehr guter Näherung:
\begin{equation}
    L(x) = \frac{N^2}{R_{m,ges}} \approx \frac{N^2 \cdot \mu_0 \cdot A_{\delta}}{\delta(x)}
\end{equation}
% -----------------------------------------------------------------------------
% 5.1.3 Mechanische Dimensionierung
\subsubsection{Mechanische Dimensionierung des Spindeltriebs und Motormoments}
\label{sec:mechanik_spindel}

Um den massiven Eisenkern (geschätzte Masse $m_{Kern} \approx 15\,\text{kg}$) stufenlos und präzise in die Kupferspule einzutauchen, wird ein linearer Spindeltrieb verwendet, der von einem Schrittmotor angetrieben wird. 

Durch Umstellen der Induktivitätsgleichung nach dem Luftspalt $\delta(x)$ lässt sich der exakte Verfahrweg berechnen. Um den Strom bei kaltem System auf $400\,\text{A}$ zu drosseln, ist eine Induktivität von ca. $20,0\,\mu\text{H}$ erforderlich, was einem minimalen Luftspalt von $\delta_{min} \approx 13,2\,\text{mm}$ entspricht. Erreicht das System seinen thermischen Beharrungszustand (630\,A Zielstrom), muss die Induktivität auf ca. $7,0\,\mu\text{H}$ sinken. Der Motor zieht das Joch folglich auf $\delta_{max} \approx 37,8\,\text{mm}$ heraus. Der Stellmotor hat somit einen hochpräzisen, nutzbaren linearen Stellweg von knapp $25\,\text{mm}$ zur Verfügung.

Für die spielfreie Positionierung wird eine Trapezgewindespindel der Nenngröße Tr 16x4 (Flankendurchmesser $d_2 = 14\,\text{mm}$, Steigung $P = 4\,\text{mm}$) spezifiziert. Die Selbsthemmung des Trapezgewindes ist ein entscheidendes Sicherheitsmerkmal: Fällt die Steuerspannung aus, verhindert die Gewindereibung, dass der schwere Kern unkontrolliert in die Spule stürzt und einen maximalen Stromsprung auslöst.
Das erforderliche Antriebsmoment $M_A$ zur Aufwärtsbewegung (gegen die Last) berechnet sich unter Berücksichtigung der Gewichtskraft ($F_G$), der Reluktanzkraft ($F_{mag}$), des Steigungswinkels $\alpha$ und des Reibungswinkels $\rho'$:
\begin{equation}
    M_A = (F_G + F_{mag}) \cdot \frac{d_2}{2} \cdot \tan(\alpha + \rho')
\end{equation}
Setzt man für $F_G \approx 147\,\text{N}$ an und geht von einer maximalen magnetischen Zugkraft von $300\,\text{N}$ aus, resultiert eine axiale Gesamtkraft von rund $450\,\text{N}$. Selbst unter diesen Extrembedingungen reicht ein industrieller Schrittmotor der Baugröße NEMA 23 oder NEMA 34 mit einem Haltemoment von $2\,\text{bis}\,4\,\text{Nm}$ aus.

% -----------------------------------------------------------------------------
% 5.2 Automatisierungskonzept mit Siemens ET200 SP
% -----------------------------------------------------------------------------
\subsection{Automatisierungskonzept mit Siemens ET200 SP}
\label{sec:automatisierung_et200sp}

Die Überführung dieser elektromechanischen Hardware in einen vollautomatisierten Prüfprozess erfolgt über das dezentrale Peripheriesystem Siemens ET200 SP. Dieses fungiert als intelligente Schnittstelle zwischen dem Starkstrom-Prüffeld und der Software-Regelebene.
% -----------------------------------------------------------------------------
\subsubsection{Elektromagnetisch verträgliche Signalverarbeitung (4-20 mA)}
In direkter Umgebung der Hauptstromschienen des Prüfstandes treten bei Dauertests bis zu $3.200\,\text{A}$ auf. Diese massiven Wechselströme erzeugen nach dem Gesetz von Biot-Savart starke magnetische Streufelder, die in ungeschirmte Signalleitungen nach dem Induktionsgesetz transiente Störspannungen einkoppeln. 
Um die Messgenauigkeit des Prüfstroms sicherzustellen, wird das Ausgangssignal der induktiven Hochstromwandler nicht als klassisches $0-10\,\text{V}$ Spannungssignal an die Steuerung übergeben. Stattdessen wird ein Messumformer eingesetzt, der den Strom-Istwert in ein analoges $4-20\,\text{mA}$ Signal konvertiert. Der elementare physikalische Vorteil dieser Stromschleife liegt im extrem niedrigen Innenwiderstand der Auswerteschaltung, wodurch kapazitiv oder induktiv eingekoppelte Störspannungen quasi kurzgeschlossen werden und das Nutzsignal nicht überlagern (Bürdenunabhängigkeit).
% -----------------------------------------------------------------------------
\subsubsection{Signalfluss und analoge Messwertverarbeitung}
Die Digitalisierung erfolgt über ein Analogeingabemodul der ET200 SP mit einer Auflösung von 16 Bit. Die Quantisierungsstufe $q$ für den Messbereich $I_{max} = 800\,\text{A}$ beträgt:
\begin{equation}
    q = \frac{800\,\text{A}}{2^{16}} \approx 12,2\,\text{mA}
\end{equation}
Diese hohe physikalische Auflösung ist zwingend erforderlich, um das enge Norm-Toleranzband der DIN EN 61439 ($0\,\%$ bis $+5\,\%$ von 630\,A) steuerungstechnisch hochpräzise abbilden zu können.

% -----------------------------------------------------------------------------
% 5.3 Softwarearchitektur und Regelungstechnik
% -----------------------------------------------------------------------------
\subsection{Softwarearchitektur und Regelungstechnik}
\label{sec:software_regelung}

Die größte Herausforderung der Regelstrecke ist die zuvor nachgewiesene vektorielle Nichtlinearität ($L \sim 1/\delta$) der R-L-Kombination.
% -----------------------------------------------------------------------------
\subsubsection{Diskrete Signalverarbeitung (PT1-Filter)}
Das TIA Portal verarbeitet die Sensorsignale zeitdiskret. Um unruhige Stellbewegungen des Tauchkerns durch Wandlerrauschen zu verhindern, wird ein digitales Verzögerungsglied 1. Ordnung (PT1-Filter) als rekursive Differenzengleichung in SCL implementiert:
\begin{equation}
    y[k] = \alpha \cdot x[k] + (1 - \alpha) \cdot y[k-1] \quad \text{mit} \quad \alpha = \frac{T_a}{T_a + T_f}
\end{equation}
Mit einer Filterzeitkonstante von $T_f = 500\,\text{ms}$ wird hochfrequentes Rauschen zuverlässig gedämpft, ohne die thermische Dynamik des Prüfstandes zu beschneiden.
% -----------------------------------------------------------------------------
\subsubsection{PID-Regler-Parametrierung und Anti-Windup}
Für die Ausregelung des thermischen Drifts wird die Anweisung \texttt{PID\_Compact} in einem deterministischen Weckalarm (OB30, $100\,\text{ms}$ Zyklus) aufgerufen. Aufgrund der enormen thermischen Trägheit des Systems wird der Regler als reiner PI-Regler konfiguriert ($T_d = 0$). 
Da die Verfahrwege des Tauchkerns mechanisch limitiert sind, ist ein Integrator-Windup-Schutz essenziell. Kann der Soll-Strom temporär nicht erreicht werden, friert die \texttt{PID\_Compact} Anweisung den I-Anteil ein, sobald die Stellgrößenbegrenzung ($0\,\%$ oder $100\,\%$) erreicht ist. Dies verhindert ein massives Überschießen bei der Rückkehr in den normalen Regelbereich.
% -----------------------------------------------------------------------------
\subsubsection{Linearisierung der Regelstrecke (Gain Scheduling)}
Da die Induktivitätsänderung bei einem großen Luftspalt deutlich geringer ausfällt als bei einem kleinen Luftspalt, variiert die Streckenverstärkung $K_S$ des Systems signifikant. Um dennoch eine konstante und stabile Reglerdynamik zu gewährleisten, wird eine arbeitspunktabhängige Adaption der Reglerverstärkung $K_p$ (Gain Scheduling) implementiert. Der $K_p$-Wert wird dynamisch über eine stückweise lineare Interpolationstabelle in Abhängigkeit der aktuellen Kernposition $x$ skaliert.

% -----------------------------------------------------------------------------
% 5.4 Sicherheits- und Überwachungskonzept
% -----------------------------------------------------------------------------
\subsection{Sicherheits- und Überwachungskonzept}
\label{sec:sicherheitskonzept}

Der unbeaufsichtigte Dauerbetrieb von Erwärmungsprüfungen über bis zu 12 Stunden erfordert ein redundantes Sicherheitskonzept zum Schutz der Anlage und des Prüfpersonals.
% -----------------------------------------------------------------------------
\subsubsection{Signalüberwachung und Sensorik (NAMUR NE 43)}
Gemäß der Industrie-Empfehlung NAMUR NE 43 wertet die ET200 SP den Signalpegel des Analogsignals kontinuierlich aus. Fällt der Strommesswert unter $3,6\,\text{mA}$, wird eine Drahtbruchstörung sicher erkannt. Die State Machine der SPS wechselt unmittelbar in den Zustand \textit{Sicherer Halt}. Dies verhindert zuverlässig, dass der Regler bei einem Drahtbruch fälschlicherweise $0\,\text{A}$ misst und den Kern unkontrolliert zur vermeintlichen Stromerhöhung verfährt, was in der Realität zu zerstörerischen Stromspitzen im Modul führen würde.
% -----------------------------------------------------------------------------
\subsubsection{Mechanischer und thermischer Eigenschutz}
Neben den softwareseitigen Begrenzungen (Soft-Limits im TIA Portal) ist der maximale lineare Stellweg der Tauchkerndrossel hardwareseitig mit elektromechanischen Endschaltern (Öffner-Prinzip) abgesichert. Ein Überfahren dieser Positionen führt zu einem sofortigen Stop des Schrittmotors (STO - Safe Torque Off).
Zusätzlich werden der Edelstahl-Festwiderstand und die Kupferwicklung der Drossel durch applizierte PT100-Sensoren thermisch überwacht. Übersteigt die Spule den thermischen Grenzwert der Isolierstoffklasse, schaltet die Steuerung das vorgeschaltete Hauptschütz des Prüfabgangs verzögerungsfrei ab.




